
\medskip
\phantomsection\label{prop:sga3-VIB-3.9}%
\noindent\textbf{Proposition 3.9.---}
\emph{Let \(G\) be an \(S\)-group locally of finite presentation.
Suppose that \(G^0\) is representable. Then the canonical morphism}
\[
  i:G^0\longrightarrow G
\]
\emph{is an open immersion; moreover, \(G^0\) is quasi-compact over
\(S\).}

Since \(G^0\) is a subfunctor of the functor \(G\), the morphism
\(i\) is a monomorphism, hence is radicial. From the definition of
the functor \(G^0\), one checks immediately from the definition
(EGA~IV\(_4\), 17.1.1) that \(i\) is formally \'etale (noting that
\(G^0\) is the image of \(i\) in \(G\)). Finally, by the
characterization (EGA~IV\(_3\), 8.14.2(c)) of \(S\)-schemes locally
of finite presentation in terms of the functors they represent, and
by \SGARefLink{cor:sga3-VIB-3.8}{3.8}, \(G^0\) is locally of finite presentation over \(S\), since
\(G\) is. Hence \(i\) is locally of finite presentation
(cf.\ EGA~IV\(_1\), 1.4.3). It is therefore a radicial \'etale
morphism, and hence an open immersion (EGA~IV\(_4\), 17.9.1).

The last assertion follows from \hyperref[cor:sga3-VIB-3.6]{3.6}.

\medskip
\phantomsection\label{thm:sga3-VIB-3.10}%
\noindent\textbf{Theorem 3.10.---}
\emph{Let \(G\) be an \(S\)-group. The following conditions are
equivalent:}
\begin{enumerate}
\item[\textup{(i)}] \emph{\(G\) is smooth over \(S\) at the points
  of the identity section.}
\item[\textup{(ii)}] \emph{\(G\) is flat and locally of finite
  presentation over \(S\) at the points of the identity section,
  and, for every \(s\in S\), \(G_s\) is smooth over \(\kappa(s)\).}
\item[\textup{\SGAInlineRefTarget{thm:sga3-VIB-3.10:item:iii}(iii)}] \emph{There is an open group subscheme \(G'\)
  of \(G\) that is smooth over \(S\).}
\item[\textup{\SGAInlineRefTarget{thm:sga3-VIB-3.10:item:iv}(iv)}] \emph{\(G^0\) is representable by an open
  subscheme of \(G\) that is smooth over \(S\).}
\end{enumerate}

It is clear that
\[
  \textup{(iv)}\Longrightarrow\textup{(iii)}
  \Longrightarrow\textup{(i)},
\]
and, by \hyperref[cor:sga3-VIB-1.3.1]{1.3.1} and \hyperref[cor:sga3-VIB-2.4]{2.4}, \textup{\SGARefLink{cor:sga3-VIB-1.3.1:item:i}{(i)}} implies both \textup{\SGARefLink{cor:sga3-VIB-1.3.1:item:ii}{(ii)}} and
\textup{(iii)}. Moreover, \textup{(ii)} implies \textup{(i)} by
EGA~IV\(_4\), 17.5.1.

It remains to prove that \textup{\SGARefLink{thm:sga3-VIB-3.10:item:iii}{(iii)}} implies \textup{\SGARefLink{thm:sga3-VIB-3.10:item:iv}{(iv)}}.
\hyperref[lem:sga3-VIB-3.10.1]{Lemma~3.10.1} below shows that \(G'\) contains \(\underline{G^0}\) and that
\[
  G'^0=G^0.
\]
It therefore suffices to show that \(\underline{G^0}\) is open in \(G\), for we
have already seen in \hyperref[case:sga3-VIB-3.4]{3.4} that \(G^0\) will then be represented by
the smooth group subscheme induced by \(G'\) on the open subset
\[
  \underline{G'^0}=\underline{G^0}.
\]
We may thus suppose that \(G'=G\).

To show that \(\underline{G^0}\) is open, it suffices to show that every
\(s\in S\) has a neighborhood \(T\) in \(S\) such that
\[
  \underline{G^0}\cap\pi^{-1}(T)
\]
is open in \(\pi^{-1}(T)\). Let \(s\in S\). Since \(G=G'\), the
morphism \(\pi\) is locally of finite presentation. As in the proof
of \SGARefLink{prop:sga3-VIB-3.7}{3.7}, we may therefore construct an open subset \(T\) of \(S\)
containing \(s\), and an open subset \(W\) of \(G\) containing
\(\epsilon(s)\), such that the induced morphism
\[
  \pi'':W\longrightarrow T
\]
is of finite presentation and admits as a section the morphism
\[
  \epsilon'':T\longrightarrow W
\]
induced by \(\epsilon\). For every \(t\in T\),
\[
  W\cap\underline{G_t^0}
\]
is then the connected component of \((\pi'')^{-1}(t)\) containing
\(\epsilon''(t)\). Since \(\pi\) is smooth, so is \(\pi''\), which
is therefore smooth and of finite presentation. By EGA~IV\(_3\),
15.6.5, the union \(W^0\) of the \(W\cap\underline{G_t^0}\), for \(t\in T\),
is open in \(W\).

On the other hand, VI\(_A\), \SGARefLink{sga3:vi:proposition:0-5}{0.5}, gives
\[
  W^0\cdot W^0=\underline{G^0}\cap\pi^{-1}(T),
\]
and we must show that this is open in \(\pi^{-1}(T)\). We may now
suppose that \(T=S\); it remains to show that \(W^0\cdot W^0\) is
open in \(G\). Since \(\pi\) is universally open, so is \(\mu\)
(VI\(_A\), \SGARefLink{sga3:VIA:0.1}{0.1}). Hence, since \(W^0\) is open in \(G\), so is
\[
  W^0\cdot W^0=\mu(W^0\times_S W^0).
\]

This result will be strengthened in \SGARefLink{cor:sga3-VIB-4.4}{4.4}.

\medskip
\phantomsection\label{lem:sga3-VIB-3.10.1}%
\noindent\textbf{Lemma 3.10.1.---}
\emph{Let \(G\) be an \(S\)-group whose fibers are locally of
finite type. Then every open group subscheme \(H\) of \(G\) contains
\(\underline{G^0}\) and satisfies}
\[
  G^0=H^0.
\]

Let \(s\in S\), and set
\[
  G'_s=H_s\cap\underline{G_s^0}.
\]
Then \(G'_s\) is an open subset of \(\underline{G_s^0}\), and is dense in
\(\underline{G_s^0}\) since \(\underline{G_s^0}\) is irreducible
(VI\(_A\), \hyperref[cor:sga3-VIB-2.4]{2.4}). Hence
\[
  G'_s\cdot G'_s=\underline{G_s^0}
\]
by VI\(_A\), \SGARefLink{sga3:vi:proposition:0-5}{0.5}, which shows that
\[
  \underline{G_s^0}=G'_s\cdot G'_s\subset H_s\cdot H_s=H_s.
\]
Thus \(\underline{G_s^0}=\underline{H_s^0}\) for every \(s\), whence
\[
  \underline{G^0}\subset H,\qquad G^0=H^0.
\]
